\section{Small inducing degrees and a projective scalar bound}
\label{sec:low-degree-projective-order}

\paragraph{Status and scope.}
These are root statements HD-LQ01--02, reserved question Q2dd.
The repaired complete proof has passed root and separate independent
whole-proof readings, including the actual inducing-base and proper-field
interfaces. The initial draft, scope repair and exact source audits are retained. The projective scalar step extends the elementary
argument of HD-TC03 without assuming a perfect normal core.
All branch entries in this section are genuine nonidentity entries
of the actual faithful inducing image. Original overgroups are not
being bounded.

\subsection*{HD-LQ01. The scalar bound needs only the projective image}

\begin{lemma}[HD-LQ01]
Let a finite subgroup \(J\le\operatorname{GL}(V)\) act irreducibly,
with \(e=\dim V\). Suppose \(T_1,\ldots,T_r\) generate \(J\),
their product is one, all are nonidentity, and
\[
 r\ge9,\qquad
 D:=\sum_i\operatorname{rank}(1-T_i)=2e+6.
 \tag{LQ.1}
\]
Put \(P=J/(J\cap\mathbb C^\times I)\) and \(s=|P|\).
Then
\[
 |J\cap\mathbb C^\times I|\mid es,
 \qquad |J|\le es^2.
 \tag{LQ.2}
\]
No primitivity, perfect-core, character-field or tensor hypothesis
is needed for this assertion.
\end{lemma}

\begin{proof}
If two entries had full moving rank, the other \(r-2\ge7\)
genuine entries would give
\(D\ge2e+7\), contradicting (LQ.1).
Thus every entry except possibly one has a fixed vector.
The projective order of every \(T_i\) divides \(s\), so
\(T_i^s\) is scalar. If \(T_i\) has a fixed vector, this scalar
is one and \(\det(T_i)^s=1\).
Taking determinants in the product-one relation gives the same
identity for the possible exceptional entry. Since the tuple
generates \(J\), its determinant image is contained in \(\mu_s\).
For any scalar \(zI\in J\), we consequently have \(z^{es}=1\).
The scalar subgroup is finite cyclic, and its order divides \(es\).
Multiplication by the projective order \(s\) proves (LQ.2).
\end{proof}

\subsection*{HD-LQ02. A uniform order ceiling in inducing degrees at most nine}

Set
\[
 B_{\rm low}=10^9,\qquad
 Q_{\rm low}=9B_{\rm low}^2=9000000000000000000.
 \tag{LQ.3}
\]

\begin{theorem}[HD-LQ02]
Let \(J\le\operatorname{GL}_e(\mathbb C)\) be the faithful
primitive inducing image of an actual rank-nine sixfold source
over the sphere, with \(e\le9\).
Then
\[
 e\ge2,\qquad |J|\le Q_{\rm low},\qquad
 r\le24,\qquad g(Y)\le1+11Q_{\rm low}
 \tag{LQ.4}
\]
for the connected inducing \(J\)-cover \(Y\to\mathbb P^1\).
The estimate includes factorwise tensor-decomposable primitive
images in this degree range and arbitrary inducing character fields.

Consequently, in HD-FQ01--03 the same integer \(Q_0\) bounds the
faithful inducing image in \emph{every} inducing degree, once the
other FQ hypotheses, including primitive factorwise tensor
indecomposability, are retained. The prime determinant exclusions
therefore apply to that entire qualified class without a lower
degree cutoff.
\end{theorem}

\begin{proof}
The actual inducing sphere source has period rank nine and therefore
at least twelve genuine branch points. This is the same source
dimension interface used in HD-TC03: a family of covers of a sphere
with \(r\) genuine marks has dimension at most \(r-3\), and its
period image cannot have greater dimension. Its Chevalley--Weil
Euler identity is \(D=2e+6\).
These are properties of the actual inducing cover, not an assignment
of twelve marks to an individual tensor factor.
If \(e=1\), the identity gives \(D=8\), whereas the twelve genuine
entries have total moving rank at least twelve. Hence \(e\ge2\).

Collins' published Theorem A gives an explicit upper bound for
\([J:Z(J)]\) for every finite primitive complex linear group.
Every entry in its degree \(2\) through \(9\) table is less than
\(10^9\); see \cite[Theorem A, p.~761]{LQ:Collins2008}.
We use this coarse envelope rather than require an optimal table
constant. Irreducibility and Schur's lemma identify \(Z(J)\)
with the scalar subgroup, so \(s=|P|\le B_{\rm low}\).
HD-LQ01 gives \(|J|\le e s^2\le9B_{\rm low}^2\).

Each genuine entry contributes at least one, so
\(r\le2e+6\le24\). If its inertia order is \(o_i\), Galois
Riemann--Hurwitz gives
\[
 2g(Y)-2=|J|\left(-2+\sum_i(1-1/o_i)\right)
       \le |J|(r-2)\le22Q_{\rm low},
 \tag{LQ.5}
\]
which proves (LQ.4).

For the all-degree FQ conclusion, first dispose of degree one
before assuming that the inducing base is a sphere.
D-MON03 excludes every monomial inducing source in the nonsplit
class under consideration. Its proof explicitly permits a proper
inducing character field \(L/K\): the actual rational isotypic
factor is \(B^{[L:K]}\), and the full generic endomorphism algebra
of the dominating cyclic genus-four component forces \(L=K\),
after which its natural polarization is split.
Thus the potentially unramified genus-four case in degree one
does not enter this nonsplit conclusion.

For \(e\ge2\), let \(h\) be the genus of the actual inducing base.
The already checked actual-cover interface D-BI-18, equations
(BI.65)--(BI.66), gives
\[
 D=2e+6-2eh,\qquad r\le D.
 \tag{LQ.6}
\]
If \(h=1\), then \(D=6\) and the pointed-curve moduli have
dimension at most six. If \(h=2\), their dimension is at most
\(3+D=9-2e\le5\). If \(h\ge3\), then \(D<0\), impossible.
The multiplicity period map factors through this actual inducing
Hurwitz space, so rank nine forces \(h=0\).
This argument uses neither equality of the original and inducing
fields nor a mark-count premise on an individual tensor factor.

In the FQ notation, \(S_0\ge10^{60}\), hence
\(A_0\ge10^{60}\), \(E_0\ge A_0\) and
\(Q_0\ge E_0>Q_{\rm low}\).
HD-FQ03 already bounds the qualified inducing images for
\(e\ge10\); (LQ.4) handles \(2\le e\le9\), with no additional
structural restriction required in these small degrees.
Together with the degree-one exclusion, the unchanged FQ bound
therefore holds in all degrees.
Applying HD-PD02 with \(Q=Q_0\) proves the unchanged prime
determinant exclusions. The extension to an arbitrary specified
polarization on a \(K\)-isogenous target still requires
\(\operatorname{End}^0=K\), exactly as in HD-PD02.
\end{proof}

\paragraph{Exact input and remaining gate.}
Collins' result is a published classification-based input, not a new
classification in this run. Root read its full Theorem A statement
and table, the definitions immediately before it, and the complete
final proof in Section5, pp.~771--774. The preceding group-theoretic
and representation estimates in Sections2--4 are cited through that
published theorem; a fresh independent reconstruction of them is not
claimed. The primary Oxford deposit is the published typeset PDF.

This removes the small-degree qualification from the FQ source-route
restriction. It does not prove that any surviving small-degree
Nielsen component has the required field descent, period rank or
detector cycles. High-degree primitive factorwise tensor-decomposable
images remain outside the all-degree FQ conclusion. No bound on the
original overgroup or arbitrary induction index is inferred.

The geometric objects remain actual finite covers, their weight-one
local systems and rational deck-correspondence projectors. The result
limits suppliers for the integral uniform condition A (cote16,
scan050; SIGN\_NEXT SN-C03/C04). Condition I (scan049; cote44,
Theorem4), the Cayley--Hamilton inverse of scans048--049 and the
existing EGA/SGA/FGA spreading interfaces do not create missing
cycles. No Markman, Severn or032 premise is introduced.
